
\ifintegriereaufgaben
\section*{Numerische Verfahren Aufgaben}\label{MA2-numerische-loesung-von-differentialgleichungen-aufgaben}
Erkl�rungen und Beispiele, die zur Bearbeitung der Aufgaben hilfreich sind, finden Sie im Skript\wennreferenzvorhanden{MA2-numerische-loesung-von-differentialgleichungen}{ (Abschnitt \abschnittsverweis{MA2-numerische-loesung-von-differentialgleichungen})} und, falls Sie nicht weiterkommen, dann schauen Sie \hyperref[MA2-numerische-loesung-von-differentialgleichungen-loesungen]{hier}.
\else
\section{Numerische Verfahren Aufgaben}\label{MA2-numerische-loesung-von-differentialgleichungen-aufgaben}

Regeln und Beispiele zu diesen Aufgaben finden Sie im Skript. Die L�sungen finden Sie \hyperref[MA2-numerische-loesung-von-differentialgleichungen-loesungen]{hier}.
\fi

\begin{exercisebox}[Explizites Euler Verfahren] %HMI 3 Blatt 4 Aufg. 4
	Berechnen Sie f�r das Anfangswertproblem
	\[
		\begin{cases}
			y(0) = 1\,,\quad & x=0\,,\\
			y' = x^2+y\,,\quad & x\not=0\,
		\end{cases}
	\]
	eine N�herungsl�sung f�r $y(2)$. Die N�herungsl�sung soll mit Hifle des Eulerschen Polygonzugverfahrens zur Schrittweite $h=1/4$ hergeleitet werden. L�sen Sie das Anfangswertproblem exakt und bestimmen Sie den relativen Fehler der Approximation f�r $x =2$.
\end{exercisebox}

\begin{exercisebox}[Verbessertes Euler Verfahren] %HMI 3 Blatt 6 Aufg. 1
	Eine Verbesserung des expliziten Euler Verfahrens lautet
	\[
		\begin{cases}
			k_1 = f(x_k, y_k)\,,\\
			k_2 = f(x_k+\frac{h}{2}, y_k + \frac{h}{2}k_1)\,,\\
			y_{k+1} =  y_k + hk_2\,.
		\end{cases}
	\]
	\begin{itemize}
		\item[(a)] Verdeutlichen Sie graphisch die Funktionsweise des Verfahrens.
		\item[(b)] Bestimmen Sie f�r das Anfangswertproblem aus der ersten Aufgabe die Approximation f�r $x=2$ zur Schrittweite \(h=1/4\) und bestimmen Sie den relativen Fehler zur exakten L�sung.
		\item[(c)] Betrachten Sie das Anfangswertproblem
	\[
		\begin{cases}
			y(0) = y_0\,,\quad & x=0\,,\\
			y' = -2at\,,\quad & x\not=0\,.
		\end{cases}
	\]
	�berzeugen Sie sich davon, dass die verbesserte Polygonzugmethode das Anfangswertproblem zu unterschiedlichen Werten von \(a = 2\) und Anfangswerten \(y_0 = \sqrt(2)\) exakt l�st (numerisch!). Was k�nnen Sie daraus schlie�en?  	 % Blatt 7 Aufg.1
	\end{itemize}
\end{exercisebox}

\begin{exercisebox}[Explizites Euler Verfahren] %HMI 3 Blatt 7 Aufg. 3
	Betrachten Sie ein explizites Euler-Verfahren f�r das Anfangswertproblem
	\[
		\begin{cases}
			y(0) = 0\,,\quad & x=0\,,\\
			y' = 4x^2-7x+3\,,\quad & x\not=0\,.
		\end{cases}
	\]
	\begin{itemize}
		\item Berechnen Sie die exakte L�sung  $\Phi$ des Anfangswertproblems.
		\item Berechnen Sie bei vorgegebener Schrittweite $h$ mit dem Euler-Verfahren N�herungswerte $y_k$ f�r $\Phi(x_k)$ und explizit die Fehler $\epsilon(x_k,h) = y_k - \Phi(x_k)$ mit $x_k = kh$.
		\item Zeigen Sie, dass der Fehler $\epsilon(x_k,h) = y_k - \Phi(x_k)$ bei festem $x$ f�r $h\to 0$ gegen Null konvergiert.
	\end{itemize}
\end{exercisebox}

\begin{exercisebox}[Implizites Euler Verfahren] %HMI 3 Blatt 7 Aufg. 1
	Betrachten Sie das Anfangswertproblem
	\[
		\begin{cases}
			y(0) = y_0\,,\quad & x=0\,,\\
			y' = -2at\,,\quad & x\not=0\,.
		\end{cases}
	\]
	Bestimmen Sie f�r das implizite Eulerverfahren den lokalen Diskretisierungsfehler
	\[
		\tau(t,h) := \frac1h \left( y(t_{k+1}) - y_{k+1} \right)
	\]
\end{exercisebox}

%\subsection{Praxis}
%\begin{exercisebox}[Moodle] % HMI III; Blatt 12, Augabe 4
%L�sen Sie das lineare Randwertproblem
%\[
%y'' - x y' +4y = x \land y(0) = 1 \land y(1) = 0 \]
%n�herungsweise mit Hilfe der Differenzenmethode, das hei�t zerlegen Sie das Intervall \([0,1]\) �quidistant mit einer Schrittweite \(h=0.2\) und approximieren Sie die auftretenden Ableitungen durch die Differentzenquotienten
%\begin{eqnarray*}
%	y'(x) &\approx& \frac{ y(x+h) - y(x-h) }{2h} \textnormal{ zentrale Differenz}\\
%	y''(x) &\approx& \frac{ y(x+h) - 2y(x) + y(x-h) }{h^2} \textnormal{ zweite Differenz}
%\end{eqnarray*}
%\end{exercisebox}

%\subsection{Praxis}
%\begin{mybox}[MC Differentialgleichungen]
%\begin{itemize}
%	\item Welcher L�sungsansatz ?
%\end{itemize}
%\end{mybox}


%https://me-lrt.de/anfangswertproblem-implizites-eulerverfahren !!!!

%
